\documentclass[twocolumn]{aastex631}
\begin{document}
\title{A residual reionization ionizing-photon-budget shortfall for star-forming galaxies at z\~{}8 (crisis systematic corner)}
\author{NebulaMind Lab (autonomous pipeline)}
\affiliation{NebulaMind Lab, \url{https://lab.nebulamind.net}}
\begin{abstract}
Reionization ionizing-photon-budget reconciliation at z\~{}8: star-forming galaxies require f\_esc=0.789 (+0.628/-0.348) to close the budget (Madau-Dickinson SFRD, log xi\_ion=25.5+/-0.15, clumping C=2-5, JWST-SFRD tail), versus indirect-proxy-inferred f\_esc=0.050 (+0.075/-0.030) from LzLCS O32/beta calibrations. Median delta(required-inferred)=+0.716 dex-frac (16-84\%: +0.361 to +1.350); 99\% of the systematic MC shows a shortfall. Conclusion: a genuine SHORTFALL remains, and the sign holds under both O32 and beta calibrations. The result is bounded by the xi\_ion x clumping x proxy-calibration systematic, not by statistics. Generated autonomously from public data (jwst) via the NebulaMind Lab runner: a bounded, reproducible, descriptive study.
\end{abstract}
\section{Introduction}
Recent studies have highlighted a potential crisis in understanding the photon budget during reionization. Muñoz et al. [Muoz2024] suggest that there may be insufficient ionizing photons produced by star-forming galaxies to account for the observed reionization, while Davies et al. [Davies2021] emphasize the increased demands on ionizing sources in absorption-dominated reionization scenarios. Park et al. [Park2022] propose a calibration method for excursion set reionization models to conserve ionizing photons more accurately. These findings underscore the importance of reconciling the ionizing photon budget during this critical period in cosmic history.
\section{Data and method}
To address this issue, we adopt a literature-anchored approach, utilizing established values from previous studies. We rely on the Madau \& Dickinson (2014) analytic fitting function for the cosmic star formation rate density (SFRD), and employ published calibrations for xi\_ion and the O32/beta f\_esc proxy [LzLCS: Chisholm+22, Flury+22; Simmonds+24]. Our method involves calculating the ionizing-photon-budget to assess whether star-forming galaxies can account for reionization at z\~{}8.
\section{Result}
Our analysis reveals that star-forming galaxies require a significantly higher escape fraction (f\_esc) of 0.789 (+0.628/-0.348) to reconcile the ionizing photon budget, assuming the Madau-Dickinson SFRD, log xi\_ion=25.5+/-0.15, and clumping factor C=2-5. In contrast, indirect-proxy-inferred f\_esc values from LzLCS O32/beta calibrations yield a much lower estimate of 0.050 (+0.075/-0.030). The median difference between the required and inferred escape fractions is +0.716 dex-frac (16-84\%: +0.361 to +1.350), with 99\% of systematic Monte Carlo simulations showing a shortfall in ionizing photons.
\begin{figure}\centering\includegraphics[width=\columnwidth]{result.png}\caption{A residual reionization ionizing-photon-budget shortfall for star-forming galaxies at z\~{}8 (crisis systematic corner)}\end{figure}
\section{Caveats}
It is essential to acknowledge the limitations of our approach. Our analysis relies on previously published values, which may introduce uncertainties due to variations in assumptions and methodologies across different studies. Additionally, our calculation depends on a single selection of parameters (xi\_ion, clumping factor) and proxy calibrations, which might not fully capture the complexity of reionization processes. Furthermore, the lack of direct observational data from surveys like JWST or SDSS may limit the accuracy of our results. These caveats highlight the need for further research to refine our understanding of the ionizing photon budget during reionization.
\section*{References}
\footnotesize
[Muoz2024] Muñoz et al. (2024). Reionization after JWST: a photon budget crisis?. bibcode 2024MNRAS.535L..37M \par [Davies2021] Davies et al. (2021). The Predicament of Absorption-dominated Reionization: Increased Demands on Ionizing Sources. bibcode 2021ApJ...918L..35D \par [Park2022] Park et al. (2022). Calibrating excursion set reionization models to approximately conserve ionizing photons. bibcode 2022MNRAS.517..192P \par [Duncan2015] Duncan et al. (2015). Powering reionization: assessing the galaxy ionizing photon budget at z < 10. bibcode 2015MNRAS.451.2030D \par [Madau2017] Madau (2017). Cosmic Reionization after Planck and before JWST: An Analytic Approach. bibcode 2017ApJ...851...50M

\end{document}
